\subsubsection{b}

Vi ønsker nå at partikkelen skal akselerere nedover som om den var på månen. Da må vi ta hensyn til forskjellen mellom tyngdeakselerasjonen på jorda og på månen.

Den totale kraften blir:
\begin{align}
F = -mg_j + mg_m
\end{align}

Setter dette inn i uttrykket for ladning:
\begin{align}
Q = \frac{F}{E} = \frac{-mg_j + mg_m}{E} = \frac{m(g_m - g_j)}{E}
\end{align}

Vi bruker verdiene:
\begin{align}
g_j = -9.81\,\text{m/s}^2, \quad g_m = -1.625\,\text{m/s}^2
\end{align}

Setter inn tallverdier:
\begin{align}
Q = \frac{0.02\,\text{kg} \cdot (-1.625 + 9.81)\,\text{m/s}^2}{100\,\text{N/C}} = 1.637 \cdot 10^{-3}\,\text{C}
\end{align}